% !TeX root = ../../Book2.tex
% 纲要标题：## 第14章　Jacobson 根、稠密性与本原环

\chapter{Jacobson 根、稠密性与本原环}
\label{b2:ch:14}
\BookChapterNumber{b2:ch:14}

\begin{chapterabstract}
Jacobson 根收集在每个简单模上都作用为零的环元素，其单位判据把这种作用性质转成可逆性问题。本章证明根在商环与矩阵环中的公式，再用非交换 Nakayama 引理提升有限生成模的生成元。随后区分本原环、半本原环与半单环，以简单模的自同态除环为正确标量建立 Jacobson 稠密性定理。代数闭域上的有限维情形进一步给出 Burnside 矩阵代数定理；具体反例说明底域、维数和量词各自的作用。
\end{chapterabstract}

\begin{learninggoals}
\begin{itemize}
\item 在极大左、右理想的交、简单模零化子与单位判据之间转换，计算商环和矩阵环的 Jacobson 根。
\item 证明并运用非交换 Nakayama 引理，分别在局部环与一般半单商上计算最少生成元个数。
\item 区分一个忠实简单模与一族简单模共同检测环元素的条件，理解半本原环的次直积表示。
\item 核对自同态除环的反环与右标量约定，使用稠密性实现有限组插值，辨明何时能得到全部矩阵代数。
\end{itemize}
\end{learninggoals}

在双数环 \(A=k[\varepsilon]/(\varepsilon^2)\) 中，\(\varepsilon\ne0\)，却在简单模 \(A/(\varepsilon)\) 上作用为零。另一方面，对任意 \(a\in A\)，\(1-a\varepsilon\) 的逆可直接写成 \(1+a\varepsilon\)。一个元素对所有简单模的作用与它参与单位判据的表现，竟然指向同一类环元素；Jacobson 根把这种现象从算例中抽出来。

\ProseParagraphBreak

对全部简单模考察作用，得到它们共同的零化子 \(J(R)\)；若改为固定一个简单模，问题就变成这个模能否忠实地反映环元素，以及环元素能在它上面实现哪些算子。前一个问题引出本原环，后一个问题引出稠密性定理。两者都从简单模出发，但“一族简单模共同忠实”与“一个简单模忠实”是不同的条件。

\ProseParagraphBreak

根也控制生成元的提升。对有限生成模 \(M\)，\(M/J(R)M\) 的任意一组生成元，无论怎样选择代表元，都会生成 \(M\)。Nakayama 引理说明为何根层不会妨碍这个提升；稠密性定理则说明为何简单模上的算子只需保持自同态除环给出的关系，就能实现任意有限组相容的指定像。

\ProseParagraphBreak

本章需要\chapref{b2:ch:12}的有限长度和局部环判据，以及\chapref{b2:ch:13}的简单模、Schur 引理和半单环结构。一般环部分不预设任何链条件；只有生成元提升、有限长度的半单性判据及有限维表示结论分别加入相应条件。单位判据中的可逆性要求“对每个环元素”成立；稠密性要求向量在自同态右除环上线性无关，不能改成较小底域上的线性无关。

\ProseParagraphBreak
本章的结论在后文有两种不同用途。\secref{b2:sec:1401}的根与单位判据及\secref{b2:sec:1402}的 Nakayama 引理支撑\chapref{b2:ch:15}的根滤过与投射覆盖，进而用于\chapref{b2:ch:16}的 Morita 理论。\secref{b2:sec:1403}的简单模和\secref{b2:sec:1404}的稠密性，尤其是自同态右除环上的有限维情形，则进入\chapref{b2:ch:17}的中心单代数和\chapref{b2:ch:18}的 Brauer 群。\secref{b2:sec:1405}的 Burnside 定理给出代数闭域上的进一步情形。

\ProseParagraphBreak
有限维代数的根可参阅 Etingof 等人的《Introduction to Representation Theory》\cite[Section 3.5, Definition 3.5.1, Proposition 3.5.2, p. 46]{EtingofEtAl2011RepresentationTheory}；代数闭域上有限维不可约表示的稠密性见\cite[Section 3.2, Corollary 3.2.1, Theorem 3.2.2(i), p. 43]{EtingofEtAl2011RepresentationTheory}。一般含幺环的 Jacobson 根与本原环可参阅 Lam 的《A First Course in Noncommutative Rings》\cite[Chapters 2 and 4]{Lam2001NoncommutativeRings}。

% !TeX root = ../../Book2.tex
% 纲要标题：### 14.1 Jacobson 根
% 纲要位置：计划纲要.md，第 1492 行。

\section{Jacobson 根}
\label{b2:sec:1401}

简单模给出环最小的非零作用对象。某个环元素可能在每个简单模上都作用为零；这些元素组成 Jacobson 根。本节把这个描述与极大左理想、可逆性及取商联系起来。环 \(R\) 均含单位元，模均幺性，不假设 \(R\) 交换或满足链条件。

% 纲要要点（供目录核对）：
%
% * 极大左理想的交及简单模作用的共同核；极大左理想商首先为模
% * 根中的可逆性判据及 Jacobson 逆公式；区分单个幂零元与理想中元素各自幂零
% * 商环、矩阵环与有限长度判据；矩阵环简单模对应、单位提升及上三角根计算
%

\subsection{极大左理想的交及等价刻画}
\label{b2:subsec:1401:01}

每个简单左模都同构于某个 \(R/L\)，其中 \(L\) 为极大左理想。\cref{b2:prop:simple-annihilator-intersection}已经证明：一个元素零化全部简单模，当且仅当它属于全部极大左理想。因而这些左理想的交正好把简单作用共同检测不到的部分收集起来。这里 \(R/L\) 首先是简单左模；\(L\) 未必为双边理想，不能一般地把它当作商环。

\begin{definition}\label{b2:def:jacobson-radical}
设 \(R\) 为含幺环。将 \(R\) 的全部极大左理想取交，所得集合称为 \(R\) 的\term[Jacobson-gen]{Jacobson 根}[Jacobson radical]，记为
\[
J(R)=\bigcap_{L\text{ 极大左理想}}L.
\]
对零环约定 \(J(0)=0\)。非零环的每个真左理想都包含于某个极大左理想。
\end{definition}
\symidx{\(J(R)\)：环 \(R\) 的 Jacobson 根}

最后的存在性可由第一卷附录A定理A.1.4的 Zorn 引理得到：包含给定真左理想的真左理想族非空，链的并仍是左理想，而且不含一，因而仍真；空链用原理想作上界。由此取到极大元。这里“极大”在所有真左理想之间比较，不要求它本身是双边理想。

\begin{proposition}\label{b2:prop:radical-simple-annihilators}
设 \(R\) 为含幺环，\(J(R)\) 为其 Jacobson 根。\(J(R)\) 等于全部幺性简单左 \(R\) 模的零化子的交，因而是双边理想。对每个这样的简单模 \(S\)，令其作用同态为
\[
\rho_S:R\longrightarrow\operatorname{End}_{\mathbb Z}(S),
\qquad \rho_S(r)(s)=rs,
\]
其中自同态取自 \(S\) 的加法群，则 \(J(R)=\bigcap_S\ker\rho_S\)。两个交都可只取简单模的同构类代表；取全部循环商 \(R/L\)、\(L\) 为极大左理想，也足以计算它们。
\end{proposition}
\begin{proof}
交的等式已在\cref{b2:prop:simple-annihilator-intersection}证明：一个元素属于所有极大左理想，恰好使它在每个简单模的每个向量上作用为零。每个简单模本来就同构于某个循环商 \(R/L\)，同构模又有相同的零化子，所以只需所述代表族。模零化子是双边理想，因为 \(xS=0\) 蕴含 \((rx)s=r(xs)=0\)、\((xr)s=x(rs)=0\)。因此这些零化子的交 \(J(R)\) 也为双边理想。作用的结合律使 \(\rho_S\) 保持乘法，且 \(\ker\rho_S=\operatorname{Ann}_R(S)\)，便得到共同核的表述。零环没有非零幺性简单模，空族的交按全环取值，等式也成立。
\end{proof}

例如，对域 \(k\) 和 \(n\ge1\)，\(R=M_n(k)\) 的列向量模 \(k^n\) 简单且忠实，所以 \(J(R)=0\)。但是 \(x=E_{12}\) 在 \(n\ge2\) 时是非零幂零元。取 \(r=E_{21}\)，就有 \(rx=E_{22}\)，它是非零幂等元，已经不再幂零。单个元素幂零，并不保证任意左倍仍幂零；这与整个理想中的元素都幂零是不同的条件。

\subsection{根中的可逆性判据}
\label{b2:subsec:1401:02}

\begin{theorem}[可逆性判据]\label{b2:thm:jacobson-unit-criterion}
设 \(R\) 为含幺环，\(J(R)\) 为其 Jacobson 根。对 \(x\in R\)，以下条件等价：
\begin{equivalences}
\item\label{b2:item:radical-member} \(x\in J(R)\)。
\item\label{b2:item:radical-unit-left} 对每个 \(r\in R\)，\(1-rx\) 为单位。
\item\label{b2:item:radical-unit-right} 对每个 \(r\in R\)，\(1-xr\) 为单位。
\end{equivalences}
此外，对任意 \(a,b\in R\)，\(1-ab\) 可逆当且仅当 \(1-ba\) 可逆；若 \(u=(1-ab)^{-1}\)，则
\[
(1-ba)^{-1}=1+bua.
\]
\end{theorem}
\begin{proof}
先证明附带的恒等事实。若 \(u=(1-ab)^{-1}\)，则 \((1-ab)u=u(1-ab)=1\) 分别给出 \(u-1-abu=0\) 与 \(u-1-uab=0\)。于是
\[
(1-ba)(1+bua)=1+b(u-1-abu)a=1,
\qquad
(1+bua)(1-ba)=1+b(u-1-uab)a=1.
\]
所以 \(1+bua\) 是 \(1-ba\) 的双侧逆；交换 \(a,b\) 便得另一方向。

对根中的元素，要先由极大左理想的定义得到左逆，再利用该左逆仍具有 \(1+J(R)\) 的形式补出右逆。设 \(j\in J(R)\)。若左理想 \(R(1-j)\) 为真，它包含在某个极大左理想 \(L\) 中，而 \(j\in L\)，导致 \(1\in L\)，矛盾。故存在 \(b\) 满足 \(b(1-j)=1\)。此时 \(b=1+bj\in1+J(R)\)。对 \(-bj\in J(R)\) 重复同一论证，有 \(c\) 使 \(cb=1\)。于是
\[
c=cb(1-j)=1-j,
\]
所以 \((1-j)b=1\)，\(1-j\) 确实可逆。若 \(x\in J(R)\)，双边理想性给出 \(rx\in J(R)\)，因而得到第二项。

若第二项成立而 \(x\notin J(R)\)，取极大左理想 \(L\) 不含 \(x\)。此时 \(L+Rx=R\)，存在 \(\ell\in L,r\in R\) 使 \(1=\ell+rx\)。所以单位 \(1-rx=\ell\) 落在真左理想中，左乘其逆即得 \(1\in L\)，矛盾。

最后将已证恒等事实用于 \(a=r,b=x\)，便得第二、三项等价。
\end{proof}

其中 \(1-ab\) 与 \(1-ba\) 的可逆性等价及逆公式通常称为 Jacobson 引理。单位判据要求 \(1-rx\) 对任意 \(r\) 都可逆，描述的是所有左倍 \(rx\) 都不破坏 \(1\) 的可逆性；它为模掉根后的单位提升提供了依据。

\begin{corollary}[根的左右对称性]\label{b2:cor:jacobson-left-right}
对任意含幺环 \(R\)，有
\[
J(R)=\bigcap_{L\text{ 极大左理想}}L
=\bigcap_{K\text{ 极大右理想}}K.
\]
\end{corollary}
\begin{proof}
极大右理想正是反环 \(R^{\mathrm{op}}\) 的极大左理想。把\cref{b2:thm:jacobson-unit-criterion}用于 \(R^{\mathrm{op}}\)，可知 \(x\) 属于全部极大右理想，当且仅当对每个 \(r\in R\)，元素 \(1-xr\) 在 \(R\) 中可逆；反环与原环的单位相同。再把同一判据用于 \(R\)，这恰是 \(x\in J(R)\) 的条件。零环按约定也满足等式。
\end{proof}

单位判据中的“对每个 \(r\)”不能省去。前面的矩阵 \(x=E_{12}\) 满足 \((1-x)^{-1}=1+x\)，但取 \(r=E_{21}\) 后，\(1-rx=1-E_{22}\) 的第二列为零，不可逆。这正是单个幂零元不必属于根的障碍。

\begin{corollary}\label{b2:cor:nil-ideal-radical}
设 \(R\) 为含幺环、\(I\subseteq R\) 为双边理想。若 \(I\) 的每个元素都幂零，则 \(I\subseteq J(R)\)。特别地，若 \(I^n=0\) 对某个正整数 \(n\) 成立，结论也成立。
\end{corollary}
\begin{proof}
对 \(x\in I,r\in R\)，理想的左乘封闭性保证 \(rx\in I\)，所以 \(rx\) 也幂零；这是单个幂零元的例子中缺少的保证。取正整数 \(m\) 使 \((rx)^m=0\)，有限几何级数给出
\[
(1-rx)^{-1}=1+rx+\cdots+(rx)^{m-1}.
\]
由可逆性判据，\(x\in J(R)\)。理想幂零时，每个元素的同一次幂都为零，故为特例。
\end{proof}

英文文献把理想中每个元素各自幂零的条件称为 nil ideal，把存在共同指数使 \(I^n=0\) 的条件称为 nilpotent ideal。本书将“幂零理想”用于后者；前者不要求各元素具有同一个幂零指数。

\ProseParagraphBreak
对域 \(k\) 和 \(n\ge1\)，在 \(k[t]/(t^n)\) 中，理想 \((\bar t)\) 幂零，故包含于根；商为域，\((\bar t)\) 又是极大理想，所以根恰为它。一般环的根未必幂零，例如 \(k[\![t]\!]\) 的单位恰为常数项非零的级数：逐次解系数可构造其逆。因此非单位元组成的理想 \((t)\) 是唯一极大理想，\(J\bigl(k[\![t]\!]\bigr)=(t)\)，但 \(t^n\ne0\) 对每个 \(n\) 成立。

\subsection{商环、矩阵环与有限长度判据}
\label{b2:subsec:1401:03}

\begin{proposition}\label{b2:prop:radical-quotient}
设 \(R\) 为含幺环、\(I\) 为其双边理想，\(\pi:R\to R/I\) 为商映射，则 \(\pi\bigl(J(R)\bigr)\subseteq J(R/I)\)。若 \(I\subseteq J(R)\)，则
\[
J(R/I)=J(R)/I.
\]
特别地，\(J\bigl(R/J(R)\bigr)=0\)。
\end{proposition}
\begin{proof}
\(R/I\) 的幺性简单左模通过 \(\pi\) 成为简单 \(R\) 模，并被 \(I\) 零化；反过来，满足 \(IS=0\) 的幺性简单 \(R\) 模 \(S\)，其作用也唯一经过 \(R/I\)。两种作用的子模相同，所以简单性不变。因此取商只保留原来简单模中被 \(I\) 零化的那些测试对象。\(J(R)\) 零化全部简单 \(R\) 模，当然也零化这些对象，由\cref{b2:prop:radical-simple-annihilators}得到包含。若 \(I\subseteq J(R)\)，则每个极大左理想都含 \(I\)；商与原像使它们与 \(R/I\) 的极大左理想一一对应。逐个取交便得等式，再取 \(I=J(R)\) 得最后断言。
\end{proof}

测试简单模的族缩小后，共同零化它们的元素可能增加，所以上述包含可以严格。例如 \(J(\mathbb Z)=0\)：若整数 \(x\ne0\)，则 \(1-3x\) 不是整数单位，单位判据排除 \(x\)。但 \(\mathbb Z/4\mathbb Z\) 的根为非零理想 \((\bar2)\)。

\begin{proposition}\label{b2:prop:matrix-radical}
设 \(R\) 为含幺环、\(n\ge1\) 为整数。每个幺性简单左 \(R\) 模 \(S\) 都给出按矩阵乘列作用的幺性简单左 \(M_n(R)\) 模 \(S^n\)。反过来，对任意幺性简单左 \(M_n(R)\) 模 \(N\)，\(E_{11}N\) 通过角环 \(E_{11}M_n(R)E_{11}\cong R\) 成为幺性简单左 \(R\) 模，且
\[
N\cong(E_{11}N)^n
\]
为左 \(M_n(R)\) 模同构。这两个构造给出两类简单模同构类的一一对应。这里 \(E_{ij}\) 为标准矩阵单位。相应地，有
\[
J\bigl(M_n(R)\bigr)=M_n\bigl(J(R)\bigr).
\]
\end{proposition}
\begin{proof}
先把两种环的简单模对应起来。取幺性简单左 \(R\) 模 \(S\)。任意非零列向量有非零坐标 \(s\)，矩阵单位可把它单独移到任意坐标，而 \(Rs=S\)，所以该向量生成全部 \(S^n\)，证明 \(S^n\) 简单。

反过来，取任意幺性简单左 \(M_n(R)\) 模 \(N\)，令 \(S=E_{11}N\)。角环 \(E_{11}M_n(R)E_{11}\) 由仅第 \((1,1)\) 个条目可能非零的矩阵组成，\(r\mapsto rE_{11}\) 是从 \(R\) 到该角环的含幺环同构，角环的单位为 \(E_{11}\)。因此按 \(r\cdot s=(rE_{11})s\) 定义的左 \(R\) 作用是幺性的。这里 \(rE_{ij}\) 表示第 \((i,j)\) 个条目为 \(r\)、其余条目为零的矩阵。

因为 \(1=\sum_iE_{ii}\)，非零 \(v\in N\) 总有某个 \(E_{ii}v\ne0\)。此时 \(E_{1i}v\ne0\)，否则再乘 \(E_{i1}\) 会得到 \(E_{ii}v=0\)，所以 \(S\ne0\)。若 \(0\ne U\subseteq S\) 是 \(R\) 子模，对任意矩阵 \(B=(b_{ji})\) 有
\[
B E_{i1}u=\sum_j E_{j1}\bigl((b_{ji}E_{11})u\bigr)
\qquad(u\in U).
\]
右端仍属于 \(\sum_jE_{j1}U\)，故这个非零子模在所有矩阵作用下稳定，由 \(N\) 简单得 \(\sum_iE_{i1}U=N\)。再施加 \(E_{11}\)，便有 \(U=S\)，所以 \(S\) 简单。

矩阵单位还给出坐标同构 \(S^n\to N\)、\((s_i)\mapsto\sum_iE_{i1}s_i\)。它的逆取 \(v\mapsto(E_{1i}v)_i\)：\(E_{1i}\) 读取第 \(i\) 个坐标，\(E_{i1}\) 将其放回原位置，由 \(E_{1j}E_{i1}=\delta_{ji}E_{11}\) 与 \(\sum_iE_{i1}E_{1i}=1\) 得到两边互逆。若 \(v=\sum_jE_{j1}s_j\)，则
\[
E_{1i}(Bv)=\sum_j(b_{ij}E_{11})s_j=\sum_j b_{ij}\cdot s_j,
\]
正是矩阵乘列的第 \(i\) 个坐标，所以这是左 \(M_n(R)\) 模同构。对原先由 \(S\) 构造的 \(S^n\)，\(E_{11}S^n\) 就是第一坐标上的 \(S\) 副本；模同构也保持这些构造，所以它们在同构类上互逆。

现在，一个矩阵属于 \(J(M_n(R))\)，当且仅当它零化全部简单模 \(S^n\)。把任意 \(s\in S\) 单独放在第 \(j\) 个坐标，就知这等价于每个条目 \(b_{ij}\) 零化全部简单 \(R\) 模。由\cref{b2:prop:radical-simple-annihilators}，又等价于所有 \(b_{ij}\in J(R)\)，得到根公式。零环的两边都为零，直接成立。
\end{proof}

环的根由极大左理想的交给出，模内全部极大子模的交也记录了所有简单商共同检测不到的部分。当这个交为零时，有限长度使这些简单商足以把整个模嵌入有限个简单模的直和，从而给出半单性；这一步中的有限性条件不能删除。

\begin{proposition}\label{b2:prop:finite-length-radical-zero}
设 \(R\) 为含幺环，\(M\) 为有限长度幺性左 \(R\) 模。若 \(M\) 的全部极大子模之交为零，则 \(M\) 半单。因此若正则左模 \({}_RR\) 有限长度，\(R/J(R)\) 为半单环。
\end{proposition}
\begin{proof}
零模直接成立。非零有限长度模有极大子模，由\cref{b2:thm:finite-length-chain-conditions,b2:prop:chain-maximal-minimal}的升链条件及极大元判据可得。为了把全部极大子模的交压缩为有限交，考察它们的所有非空有限交。降链条件使这一族中有一个极小者 \(K=\bigcap_{i=1}^r L_i\)。若 \(K\ne0\)，取 \(0\ne x\in K\)。全部极大子模交为零，故某个极大子模 \(L\) 不含 \(x\)，于是 \(K\cap L\subsetneq K\)，仍是一个有限交，矛盾。因此已有有限个极大子模满足 \(\bigcap_{i=1}^rL_i=0\)。

映射 \(M\to\bigoplus_{i=1}^rM/L_i\)、\(m\mapsto(m+L_i)_i\) 的核为 \(K=0\)，所以它单射。每个商模简单，目标半单，\cref{b2:prop:semisimple-subquotients}保证其子模也半单，故 \(M\) 半单。

令 \(B=R/J(R)\)。它作为左 \(R\) 模具有有限长度，因为有限长度对取商保持。其 \(R\) 作用已经经过 \(B\)，所以 \(R\) 子模与左 \(B\) 子模相同；这些子模在正则 \(B\) 模中又正是左理想。\cref{b2:prop:radical-quotient}给出 \(J(B)=0\)，故上述判据使它作为 \(R\) 模半单。同样因为两种标量作用的子模相同，简单性及直和分解也相同，\({}_BB\) 半单，亦即 \(B\) 是半单环。
\end{proof}

\begin{corollary}[单位的提升]\label{b2:cor:units-lift-radical}
设 \(R\) 为含幺环，\(I\) 为满足 \(I\subseteq J(R)\) 的双边理想。对任意 \(a\in R\)，\(a\) 是 \(R\) 的单位，当且仅当 \(a+I\) 是 \(R/I\) 的单位。更一般地，对任意整数 \(n\ge1\) 和矩阵 \(A\in M_n(R)\)，\(A\) 可逆，当且仅当逐条目取商所得矩阵 \(\overline A\in M_n(R/I)\) 可逆。
\end{corollary}
\begin{proof}
保持单位元的商映射将单位送到单位，所以只需证明反向。若 \(a+I\) 可逆，取其逆的一个提升 \(b\in R\)，有 \(ab=1+i\)、\(ba=1+j\)，其中 \(i,j\in I\)。由\cref{b2:thm:jacobson-unit-criterion}，\(1+i\) 与 \(1+j\) 都可逆。于是
\[
a\bigl(b(1+i)^{-1}\bigr)=1,
\qquad
\bigl((1+j)^{-1}b\bigr)a=1.
\]
一个左逆与一个右逆必相等：若 \(da=1\)、\(ac=1\)，则 \(d=d(ac)=(da)c=c\)。因此 \(a\) 有双侧逆。

逐条目取商的环同态 \(M_n(R)\to M_n(R/I)\) 满射，核为 \(M_n(I)\)，所以目标同构于 \(M_n(R)/M_n(I)\)。又由\cref{b2:prop:matrix-radical}，有 \(M_n(I)\subseteq M_n(J(R))=J(M_n(R))\)。将刚证的单位提升结论用于这个矩阵环及其理想，便得到矩阵版本。
\end{proof}

\begin{proposition}\label{b2:prop:triangular-radical}
设 \(k\) 为域、\(n\ge1\) 为整数，\(T_n(k)\) 为 \(n\times n\) 上三角矩阵组成的含幺 \(k\) 代数，\(N\) 为其中严格上三角矩阵组成的双边理想。则
\[
J(T_n(k))=N,\qquad N^n=0,\qquad T_n(k)/N\cong k^n.
\]
对 \(n\ge2\)，有 \(N^{n-1}\ne0\)，故 \(N\) 的最小幂零指数恰为 \(n\)；\(n=1\) 时 \(N=0\)。
\end{proposition}
\begin{proof}
取对角条目的映射 \(\delta:T_n(k)\to k^n\) 保持加法、乘法和单位元：上三角矩阵乘积的第 \((i,i)\) 个条目恰为两个第 \((i,i)\) 个条目的乘积。它满射，核正是 \(N\)，因此 \(N\) 是双边理想，且 \(T_n(k)/N\cong k^n\)。

严格上三角矩阵中非零条目只能沿严格递增的行列指标出现。\(n\) 个这样的矩阵相乘时，任一可能非零项都需要 \(n+1\) 个严格递增的指标，而指标总共只有 \(n\) 个，所以乘积为零，\(N^n=0\)。由\cref{b2:cor:nil-ideal-radical}，\(N\subseteq J(T_n(k))\)。\(k^n\) 的各坐标模 \(k\) 都简单，而且共同忠实，所以由\cref{b2:prop:radical-simple-annihilators}，\(J(k^n)=0\)。再用\cref{b2:prop:radical-quotient}，根在对角商中的像为零，故 \(J(T_n(k))\subseteq N\)，得到等式。

若 \(n\ge2\)，标准矩阵单位的乘积 \(E_{12}E_{23}\cdots E_{n-1,n}=E_{1n}\ne0\)，属于 \(N^{n-1}\)，所以没有小于 \(n\) 的幂零指数。\(n=1\) 时严格上三角部分为空，\(N=0\)，所述结论直接成立。
\end{proof}

% !TeX root = ../../Book2.tex
% 纲要标题：### 14.2 非交换 Nakayama 引理与生成元提升
% 纲要位置：计划纲要.md，第 1498 行。

\section{非交换 Nakayama 引理与生成元提升}
\label{b2:sec:1402}

模去 Jacobson 根后，可以先在剩余模中考察生成问题。Nakayama 引理把商模的生成性提升回来；它不要求环交换，但对一般的根理想需要模有限生成。

% 纲要要点（供目录核对）：
%
% * 对有限生成左 \(R\)-模 \(M\) 和双边理想 \(I\subseteq J(R)\)，证明 \(IM=M\Rightarrow M=0\)；整理系数核对右乘封闭性，删生成元用 \(1-a\) 的可逆性；理想幂零时无须有限生成
% * 证明 \(M=N+IM\Rightarrow M=N\) 的任意生成元提升形式；不要求子模有限生成，并区分满射判据与单射
% * 局部环中最少生成组与不可删减生成组的关系；附加半单商假设下按矩阵块容量计算最少生成元数
% * 第三卷第3.2节保留交换环行列式技巧及局部版本的完整处理；第二卷第15章只调用本节已证的非交换形式
%

\subsection{非交换 Nakayama 引理}
\label{b2:subsec:1402:01}

\begin{theorem}[非交换 Nakayama 引理]\label{b2:thm:noncommutative-nakayama}
设 \(R\) 为含幺环、\(M\) 为幺左 \(R\) 模，\(I\) 为双边理想。若 \(M\) 有限生成、\(I\subseteq J(R)\) 且 \(IM=M\)，则 \(M=0\)。若 \(I^r=0\) 对某个正整数 \(r\) 成立，则 \(IM=M\) 同样蕴含 \(M=0\)，不需要 \(M\) 有限生成。
\end{theorem}

这称为\term[Nakayama-yinli]{Nakayama 引理}[Nakayama's lemma]。\(IM\) 表示有限和 \(\sum_j a_jm_j\)、\(a_j\in I\) 构成的集合。它是左子模，因为对 \(c\in R\)，有 \(c\sum_j a_jm_j=\sum_j(ca_j)m_j\)，而 \(I\) 对左乘封闭。有限生成情形的证明要从一组最少生成元中删去一个元素；\(1-a\) 的可逆性正用来解出这个元素。

\begin{proof}
先设 \(M\) 有限生成且 \(I\subseteq J(R)\)。若 \(M\ne0\)，其有限生成组的元素个数构成非空自然数集，故可由自然数的良序性选出元素个数最少的一组 \(x_1,\ldots,x_n\)，此时 \(n\ge1\)，无须另加 Noether 条件。由 \(x_n\in IM\)，写成有限和 \(x_n=\sum_jb_jm_j\)，其中 \(b_j\in I\)，再将各 \(m_j\) 展开为 \(m_j=\sum_i r_{ji}x_i\)。于是
\[
x_n=\sum_{i=1}^n a_i x_i,\qquad
a_i=\sum_j b_jr_{ji}\in I.
\]
这里系数 \(b_j\) 位于 \(r_{ji}\) 左侧，故 \(a_i\in I\) 用的是 \(I\) 对右乘封闭。所以 \((1-a_n)x_n=\sum_{i<n}a_ix_i\)。\cref{b2:thm:jacobson-unit-criterion}保证 \(1-a_n\) 为单位，左乘其逆便把 \(x_n\) 写成其余生成元的线性组合，与最少性矛盾。\(n=1\) 时右侧为空和，这同样表明唯一生成元为零，矛盾。

若 \(I^r=0\)，则由有限和的定义及作用的结合律有 \(I(IM)=I^2M\)。迭代 \(M=IM\) 得
\[
M=IM=I^2M=\cdots=I^rM=0.
\]
这个论证不涉及生成组，因而适用于任意幺左模。
\end{proof}

对于一般的 \(I\subseteq J(R)\)，有限生成条件不能省去。设 \(k\) 为域、\(t\) 为不定元，取 \(R=k[t]_{(t)}\)，即分母在 \(t=0\) 处不为零的有理函数组成的局部环。其非单位恰为理想 \((t)\)，故 \(J(R)=(t)\)。作为 \(R\) 模，非零的 \(M=k(t)\) 却满足 \(tM=M\)。

\ProseParagraphBreak
它不有限生成。把非零有理函数写成 \(f=t^a g(t)/h(t)\)，其中 \(a\in\mathbb Z\)、\(g,h\in k[t]\) 且 \(g(0)h(0)\ne0\)，定义它在 \(t=0\) 处的阶为 \(v_t(f)=a\)，并取 \(v_t(0)=+\infty\)。乘积的阶相加，和的阶不低于各项阶的最小值：提出共同的最低次 \(t\) 幂后，加法抵消只可能使阶升高。\(R\) 的元素均有非负阶，所以若有限组 \(f_1,\ldots,f_n\) 的阶均不低于 \(-N\)，其中 \(N\ge0\)，则用 \(R\) 中元素作系数相乘并相加仍不能使阶低于 \(-N\)。它们因而不能生成 \(t^{-N-1}\)。

\subsection{生成元提升与有限生成商模}
\label{b2:subsec:1402:02}

\begin{corollary}[生成元提升]\label{b2:cor:nakayama-generators}
设 \(R\) 为含幺环、\(M\) 为幺左 \(R\) 模、\(N\subseteq M\) 为子模，\(I\) 为双边理想。假设 \(M\) 有限生成且 \(I\subseteq J(R)\)，或者 \(I^r=0\) 对某个正整数 \(r\) 成立。若 \(M=N+IM\)，则 \(M=N\)。因此，在上述任一假设下，\(M/IM\) 的一组生成元在 \(M\) 中的任意提升都生成 \(M\)；幂零理想的情形不要求 \(M\) 或该生成组有限。
\end{corollary}
\begin{proof}
在 \(M\) 有限生成且 \(I\subseteq J(R)\) 的情形，商模 \(M/N\) 由原有限生成组的像生成，所以仍有限生成，不需要 \(N\) 有限生成。等式 \(M=N+IM\) 给出 \(I(M/N)=M/N\)，由\cref{b2:thm:noncommutative-nakayama}，商模为零。若 \(I^r=0\)，则同一定理的幂零版本直接用于 \(M/N\)，也得商模为零，无须有限生成。

对于后一断言，任意选取所给商模生成元的提升 \(x_\lambda\)，并令 \(N=\sum_\lambda Rx_\lambda\)。每个 \(m+IM\) 都是这些商类的有限线性组合，故 \(m\) 减去相应提升的线性组合落在 \(IM\) 中。这就给出 \(M=N+IM\)，应用前一断言即可；整个论证没有要求提升具有额外性质。
\end{proof}

\begin{proposition}\label{b2:prop:nakayama-surjectivity-test}
设 \(R\) 为含幺环，\(M\) 为有限生成幺左 \(R\) 模、\(P\) 为任意幺左 \(R\) 模，\(f:P\to M\) 为 \(R\) 模同态，\(I\subseteq J(R)\) 为双边理想，则 \(f\) 满射当且仅当诱导的 \(P/IP\to M/IM\) 满射。
\end{proposition}
\begin{proof}
正向由代表元直接提升。反向，商映射满射意味着 \(M=f(P)+IM\)。对有限生成目标 \(M\) 应用\cref{b2:cor:nakayama-generators}，得 \(f(P)=M\)。
\end{proof}

这一判据不能用来检测单射。设 \(k\) 为域，取 \(R=k[t]/(t^2)\)、\(I=J(R)=(\bar t)\)、\(P=R\)、\(M=R/I\)，令 \(f:P\to M\) 为自然商映射。诱导映射 \(P/IP\to M/IM\) 是 \(k\to k\) 的恒等映射，但 \(\ker f=I\ne0\)。核恰好位于来源模的 \(IP\) 中，模去这一层后便完全消失，所以剩余映射的单射性不能恢复原映射的单射性。

\subsection{剩余模与极小生成集}
\label{b2:subsec:1402:03}

沿用\cref{b2:def:local-ring}的约定，局部环是非零环，且其非单位组成一个双边理想 \(\mathfrak m\)。此时 \(R/\mathfrak m\) 是除环：每个非零商类都有单位代表元。任意真左理想不含单位，所以包含于 \(\mathfrak m\)；后者已极大，因而是唯一极大左理想，\(J(R)=\mathfrak m\)。

\ProseParagraphBreak
对任意含幺环 \(R\) 及有限生成幺左 \(R\) 模 \(M\)，记 \(\mu_R(M)\) 为生成 \(M\) 所需的最少元素个数，约定 \(\mu_R(0)=0\)。元素数等于 \(\mu_R(M)\) 的生成组称为最少生成组；删去任一元素就不再生成的组称为不可删减生成组。局部环中，剩余模是除环上的向量空间，可以先求它的一组基，再由 Nakayama 引理提升回原模。
\symidx{\(\mu_R(M)\)：有限生成模所需的最少生成元个数}

\begin{proposition}\label{b2:prop:local-minimal-generators}
设 \(R\) 为非零含幺局部环、\(D=R/J(R)\) 为其剩余除环，\(M\) 为有限生成幺左 \(R\) 模。剩余模 \(M/J(R)M\) 是左 \(D\) 向量空间，以下 \(\dim_D\) 表示左 \(D\) 维数。生成 \(M\) 所需的最少元素个数为
\[
\mu_R(M)=\dim_D M/J(R)M.
\]
在上述局部环条件下，生成组不可删减，当且仅当其像为剩余模的一组左 \(D\) 基。因此不可删减生成组与最少生成组相同。
\end{proposition}
\begin{proof}
任意生成组的像都生成剩余模，所以元素数不少于右侧的维数。反向，提升一组有限左 \(D\) 基，由\cref{b2:cor:nakayama-generators}得到 \(M\) 的生成组，故最小数恰为该维数。若生成组的像线性相关，取一个有限非平凡关系 \(\sum_i a_i\bar x_i=0\)，其中 \(a_i\in D\)，并选 \(a_j\ne0\)，则
\[
\bar x_j=-a_j^{-1}\sum_{i\ne j}a_i\bar x_i.
\]
由于这是左 \(D\) 向量空间，逆元须乘在左侧。删去 \(x_j\) 后，其余元素仍生成剩余模，再用提升推论便仍生成 \(M\)，所以原组可再删减。若像为基，删去任意一项连剩余模都不能生成，原组因而不可删减。
\end{proof}

一般环中不可删减生成组未必最少。设 \(k\) 为域，取 \(R=k\times k\) 的正则左模。\((1,0),(0,1)\) 生成 \(R\)，删去任一个都会失去一个坐标方向，所以不可删减；但 \((1,1)\) 是单位，单独就生成 \(R\)。因此 \(\mu_R(R)=1\)，这组不可删减生成组却有两个元素。

\ProseParagraphBreak
对于一般环，\cref{b2:prop:radical-quotient}只保证 \(J(R/J(R))=0\)，并不保证 \(R/J(R)\) 半单。例如 \(J(\mathbb Z)=0\)，而 \(\mathbb Z\) 的正则模不是半单模：它不含非零简单子模，因为每个非零子模 \(n\mathbb Z\) 都含有真非零子模 \(2n\mathbb Z\)。下面的矩阵块公式因而需要附加半单商的假设。有限维代数的正则模有有限长度，由\cref{b2:prop:finite-length-radical-zero,b2:thm:artin-wedderburn}，其半单商才自动具有这种有限矩阵块分解。

\ProseParagraphBreak
先看 \(R=M_2(k)\)、\(S=k^2\) 及 \(M=S^3\)，其中 \(k\) 为域，\(R\) 同时左乘三个列向量。任何单个元素 \(v\) 生成的子模都是 \(R\to M\)、\(A\mapsto Av\) 的像，其 \(k\) 维数至多为 \(\dim_kR=4\)，而 \(\dim_kM=6\)，所以一个元素不能生成 \(M\)。另一方面，取标准列向量 \(e_1,e_2\)，则
\[
v=(e_1,e_2,0),\qquad w=(0,0,e_1)
\]
生成 \(M\)：\(Av=(Ae_1,Ae_2,0)\) 的前两个分量是 \(A\) 的两列，可以独立指定；\(Bw=(0,0,Be_1)\) 可以任意指定第三个分量。因此 \(\mu_R(S^3)=2\)，三个简单直和分量并不意味着需要三个模生成元。

\ProseParagraphBreak
这个现象来自正则模所含的简单模重数。若 \(B=\prod_{i=1}^sM_{n_i}(D_i)\) 为有限个矩阵块的直积，则一份正则左模 \(B\) 包含 \(n_i\) 份第 \(i\) 类简单列模 \(S_i\)。一个生成元给出一份 \(B\) 到目标模的同态，因而至多提供 \(n_i\) 份这种简单模。若目标含 \(m_i\) 份 \(S_i\)，生成元数至少为 \(\lceil m_i/n_i\rceil\)。同一组生成元要同时满足所有类型的重数要求，所以应取这些下界的最大值；下面的证明还会说明这个最大值确实足够。

\begin{proposition}\label{b2:prop:semisimple-top-generators}
设 \(R\) 为非零含幺环，\(M\) 为有限生成幺左 \(R\) 模。设正整数 \(s,n_1,\ldots,n_s\)、除环 \(D_1,\ldots,D_s\) 及非负整数 \(m_1,\ldots,m_s\) 满足
\[
R/J(R)\cong\prod_{i=1}^s M_{n_i}(D_i),\qquad
M/J(R)M\cong\bigoplus_{i=1}^s S_i^{m_i},
\]
其中 \(S_i=D_i^{n_i}\) 为相应矩阵块的简单列模，则
\[
\mu_R(M)=\max_i\left\lceil\frac{m_i}{n_i}\right\rceil.
\]
零模的最少生成元个数取零。
\end{proposition}
\begin{proof}
记半单商为 \(B\)。由\cref{b2:thm:artin-wedderburn}的正则模分解，\({}_BB\cong\bigoplus_i S_i^{n_i}\)，故 \(B^r\cong\bigoplus_i S_i^{rn_i}\)。\cref{b2:prop:semisimple-multiplicity}给出的简单模重数唯一性说明，\(\bigoplus_iS_i^{m_i}\) 是 \(B^r\) 的商，当且仅当每个 \(m_i\le rn_i\)：充分性取对应分量投影，必要性由满射的分裂及直和重数相加得到。这正是剩余模可由 \(r\) 个元素生成的条件。原模生成元映到剩余模生成元，而由\cref{b2:cor:nakayama-generators}，后者的提升又生成原模，所以两者的最少数相同，得到公式。
\end{proof}

交换局部环的常用 Nakayama 引理已经包含在以上结论中。交换性还允许另一种证明方法：用有限生成元的关系矩阵及其伴随矩阵，把若干线性关系转成一个同时零化全部生成元的标量等式。\cref{b2:prop:adjugate-identity}提供相应的矩阵恒等式；第三卷第3.2节将把这种行列式技巧及其整性应用完整展开。

\ProseParagraphBreak
对于有限维非交换代数，下一章先使根的幂稳定，再把稳定项看成有限生成左模，用\cref{b2:thm:noncommutative-nakayama}证明该项为零。

% !TeX root = ../../Book2.tex
% 纲要标题：### 14.3 本原环与半本原环
% 纲要位置：计划纲要.md，第 1505 行。

\section{本原环与半本原环}
\label{b2:sec:1403}

模去 Jacobson 根以后，每个非零环元素至少能被某个简单模检测到。更强的问题是，能否用同一个简单模同时检测全部元素？本原环要求后者，半本原环只要求前者。它们与半单环的区别，将通过具体作用和无限乘积表现出来。

% 纲要要点（供目录核对）：
%
% * 忠实简单模与本原环；区分单个忠实作用、一族共同忠实作用及半单性
% * 半本原环的次直积结构；取极大左理想中最大的双边理想形成真正的本原商环
% * 半本原但非半单的实例；有限坐标满射与无限直积满射的区别
%

\subsection{忠实简单模与本原环}
\label{b2:subsec:1403:01}

\begin{definition}\label{b2:def:primitive-semiprimitive}
设 \(R\) 为非零含幺环。若 \(R\) 有忠实的幺性简单左模，则称它为\term[zuo-benyuan-huan]{左本原环}[left primitive ring]。若 \(J(R)=0\)，则称它为\term[banbenyuan-huan]{半本原环}[semiprimitive ring]。本节简称“本原”时均指左本原；右本原是另一条件，一般不能自动互换\footnote{右本原而非左本原的例子见\cite[pp. 473--475]{Bergman1964PrimitiveOneSided}；对该环取反环得到相反方向的例子。}。
\end{definition}

由\cref{b2:prop:radical-simple-annihilators}，半本原性等价于：对每个 \(0\ne r\in R\)，都存在一个幺性简单左模 \(S\) 和其中的向量 \(s\)，使 \(rs\ne0\)。这里 \(S\) 可以随 \(r\) 改变。本原性则要求先找到同一个简单模 \(S\)，使对每个 \(0\ne r\in R\)，都能在这个 \(S\) 中找到 \(s\) 满足 \(rs\ne0\)。因此每个本原环都半本原。对任意简单左模 \(S\)，商环 \(R/\operatorname{Ann}_R(S)\) 在 \(S\) 上的作用忠实且简单，所以它本原。

\ProseParagraphBreak
这里的“本原环”与下一章\cref{b2:def:primitive-idempotent}中的“本原幂等元”有不同的定义。前者要求一个忠实简单模，后者要求一个非零幂等元不能继续分成两个非零正交幂等元之和，不能仅凭同一个名称互换条件。

\begin{proposition}\label{b2:prop:commutative-primitive-field}
非零交换含幺环 \(R\) 左本原，当且仅当 \(R\) 是域。
\end{proposition}
\begin{proof}
若 \(S\) 忠实且简单，取 \(0\ne s\in S\)，满射 \(R\to S\)、\(r\mapsto rs\) 的核为极大理想 \(\mathfrak m\)，所以 \(S\cong R/\mathfrak m\)。交换性使这个商模的零化子恰为 \(\mathfrak m\)，忠实性便给出 \(\mathfrak m=0\)，故 \(R\) 是域。反向，一个域在自身上的正则模简单且忠实。
\end{proof}

在非交换情形，对任意除环 \(D\) 和整数 \(n\ge1\)，矩阵环 \(M_n(D)\) 都本原：列模 \(D^n\) 简单且忠实，简单性由矩阵单位移动非零坐标并乘除环元素得到。它在 \(n>1\) 时不是除环，说明上一个命题的交换性不能去掉。

\ProseParagraphBreak
非零半单环都半本原：左正则模是简单模的直和，而 \(J(R)\) 零化每个简单分量，所以 \(J(R)R=0\)，取右侧元素为 \(1\) 得 \(J(R)=0\)。但半单环未必本原：对任意域 \(k\)，环 \(k\times k\) 是两个域的乘积，因而半单；它是交换环却不是域，由上一个命题可知它不本原。

\subsection{半本原环的次直积结构}
\label{b2:subsec:1403:02}

\begin{definition}\label{b2:def:subdirect-product}
设 \(I\) 为指标集，\(R\) 与 \(R_i\)、\(i\in I\) 为含幺环，\(\phi:R\to\prod_{i\in I}R_i\) 为保幺环同态。若 \(\phi\) 单射，且它与每个因子投影的复合都满射，则称 \(\phi\) 把 \(R\) 表示为这些环的一个\term[cizhiji]{次直积}[subdirect product]。这里不要求 \(\phi\) 到整个直积满射。
\end{definition}

单射使像成为直积的一个子环，坐标满射使每个单独的因子都能取遍；不同坐标仍可能受到兼容条件限制。例如对任意域 \(k\)，对角嵌入 \(k\to k\times k\)、\(a\mapsto(a,a)\) 的两个坐标映射均满射，但它的像只包含两个坐标相等的元素。

\begin{theorem}\label{b2:thm:semiprimitive-subdirect}
非零含幺环 \(R\) 半本原，当且仅当它可以表示为一族左本原环的次直积。
\end{theorem}
\begin{proof}
对每个极大左理想 \(L\)，\(R/L\) 是简单左模，但 \(L\) 未必是双边理想，不能直接把它当作商环。为了得到一个在这个简单模上忠实作用的环商，应当模去作用的核，于是令
\[
I_L=\operatorname{Ann}_R(R/L)
=\{a\in R:aR\subseteq L\}.
\]
这个集合对加减封闭；若 \(aR\subseteq L\)，则对 \(r,x\in R\)，有 \((ra)x=r(ax)\in L\) 以及 \((ar)x=a(rx)\in L\)，故 \(I_L\) 是双边理想。取 \(x=1\) 可知 \(I_L\subseteq L\)。若双边理想 \(I\subseteq L\)，则每个 \(a\in I\) 都满足 \(aR\subseteq I\subseteq L\)，所以 \(I\subseteq I_L\)。这证明 \(I_L\) 是 \(L\) 中最大的双边理想。

商环 \(R/I_L\) 在 \(R/L\) 上的作用忠实，因为模去的恰是原作用的核；商前后子模相同，所以作用仍简单，故 \(R/I_L\) 本原。全部商映射合成
\[
\phi:R\longrightarrow\prod_L R/I_L.
\]
每个坐标映射满射，整体的核是 \(\bigcap_LI_L=J(R)\)，由\cref{b2:prop:radical-simple-annihilators}已知。故 \(J(R)=0\) 时它就是次直积表示。极大左理想组成集合，因此这里的乘积是集合编号的乘积。

反过来，设 \(\phi:R\to\prod_iR_i\) 为次直积表示，各 \(R_i\) 本原且半本原。每个坐标满射把 \(J(R)\) 映进 \(J(R_i)=0\)，这是\cref{b2:prop:radical-quotient}对其核所给的包含。于是 \(\phi\bigl(J(R)\bigr)=0\)，而 \(\phi\) 单射，故 \(J(R)=0\)。
\end{proof}

整数的模素数映射给出一个同时满足所有有限坐标要求、却不能任意指定无限坐标的例子：
\[
\mathbb Z\longrightarrow\prod_{p\text{ 素数}}\mathbb F_p.
\]
它单射，因为一个非零整数不可能被所有素数整除；各坐标显然满射。对任意有限多个不同素数，第一卷定理10.4.3的中国剩余定理允许同时任意指定这些坐标。但一个无限坐标族未必来自同一个整数：将奇素数坐标全取零、模二坐标取一就不可能提升。任何可能提升的整数须被所有奇素数整除，只能为零，又与模二条件矛盾。因此有限组同余条件的可解性，不能直接推广成对整个无限直积的满射性。

\subsection{非半单的半本原环}
\label{b2:subsec:1403:03}

\(\mathbb Z\) 的根为零，已在\secref{b2:sec:1401}由单位判据算出；它的正则模不半单，也已在\secref{b2:sec:1402}用不存在非零简单子模说明。另一方面，它不是域，由\cref{b2:prop:commutative-primitive-field}可知也不本原。所以这个半本原环既不本原，也不半单。

\ProseParagraphBreak
本原也不蕴含半单。令 \(k\) 为域，\(V\) 为有可数基 \(e_1,e_2,\ldots\) 的 \(k\) 向量空间，取 \(R=\operatorname{End}_k(V)\)。任意 \(0\ne v\in V\) 都可扩充为基；把 \(v\) 送到任意给定的 \(w\)，把这组基中的其余向量送到零，就唯一确定一个线性算子。无限维时扩充基使用第一卷附录A定理A.1.4的 Zorn 引理，和构造一般向量空间的基相同。故 \(Rv=V\)，自然模 \(V\) 简单；作用按定义忠实，所以 \(R\) 本原。

\ProseParagraphBreak
要检测这个环是否左 Artin，可以要求算子杀掉越来越多的指定基向量。这样的条件给出左理想
\[
L_n=\bigl\{\,T\in R:T(e_1)=\cdots=T(e_n)=0\,\bigr\}
\]
条件每增加一个，允许的算子就减少，形成无限严格降链 \(R\supsetneq L_1\supsetneq L_2\supsetneq\cdots\)。这些集合对加减和左乘稳定，因为复合 \(ST\) 仍将指定的基向量送到零；\(L_n\supsetneq L_{n+1}\) 的严格性可由只将 \(e_{n+1}\) 送到 \(e_1\)、将其余基向量送到零的算子核对，\(R\supsetneq L_1\) 则由恒等算子核对。因此 \(R\) 不是左 Artin 环，而\cref{b2:thm:semisimple-ring-characterizations}中的半单环均为左 Artin 环。这个例子说明本原性也不蕴含半单性，并且忠实简单模可能无限维。这里在有限组向量上指定算子的行为，正是下一节有限拓扑所记录的条件。

% !TeX root = ../../Book2.tex
% 纲要标题：### 14.4 Jacobson 稠密性定理
% 纲要位置：计划纲要.md，第 1229 行。

\section{Jacobson 稠密性定理}
\label{b2:sec:1404}

在一个简单模上，环元素可以把任意非零向量送到任何向量。稠密性定理把这个单向量事实推广为有限组向量的同时插值。能够自由指定的向量组须对自同态除环线性无关，而不只是对某个较小的底域无关。

% 纲要要点（供目录核对）：
%
% * 简单模上的作用与自同态除环
% * 稠密性定理的代数形式
% * 有限维情形的矩阵代数结论
%

\subsection{简单模的作用与自同态除环}
\label{b2:subsec:1404:01}

设 \(S\) 为简单左 \(R\) 模。由\cref{b2:lem:schur}，\(\Delta=\operatorname{End}_R(S)\) 是除环。为了把自同态看作右标量，规定
\[
D=\Delta^{\mathrm{op}},\qquad v\cdot d=f_d(v),
\]
其中 \(f_d\) 是 \(d\) 对应的 \(R\) 模自同态。因为 \(f_{d_1d_2}=f_{d_2}f_{d_1}\)，这给出右 \(D\) 向量空间结构。左 \(R\) 作用与它交换，所以作用同态的目标可缩为
\[
\rho:R\longrightarrow\operatorname{End}_D(S_D).
\]
这里 \(\operatorname{End}_D\) 表示右 \(D\) 线性自同态。
\symidx{\(D=\operatorname{End}_R(S)^{\mathrm{op}}\)：简单左模的右标量除环}

一个有限向量组 \(v_1,\ldots,v_n\) 右 \(D\) 线性无关，意指 \(\sum_i v_i d_i=0\) 只可能在全部 \(d_i=0\) 时成立。若它们原本有关系，任何 \(r\in R\) 都必保留该关系，故不能任意指定它们的像。

\ProseParagraphBreak
例如，令 \(R=\mathbb C\) 在 \(S=\mathbb C\) 上左乘，并同时把 \(S\) 看成二维实空间。向量 \(1,i\) 对实数线性无关，却满足右复关系 \(1\cdot i-i\cdot1=0\)。不可能有同一个左乘把 \(1\) 送到零而把 \(i\) 送到一。这就是必须使用正确标量除环的原因。

\subsection{稠密性定理的代数形式}
\label{b2:subsec:1404:02}

\begin{definition}\label{b2:def:finite-topology}
设 \(D\) 为除环、\(S\) 为右 \(D\) 向量空间。对 \(T\in\operatorname{End}_D(S_D)\) 和有限组 \(v_1,\ldots,v_m\in S\)，令
\[
U(T;v_1,\ldots,v_m)
=\bigl\{\,F\in\operatorname{End}_D(S_D):F(v_i)=T(v_i)\text{ 对所有 }i\,\bigr\}.
\]
以这些集合为各 \(T\) 的基本邻域，所得拓扑称为\term[youxian-tuopu]{有限拓扑}[finite topology]。有限个这样的条件可以合并成一组向量条件，因此确实给出邻域基。
\end{definition}

\begin{theorem}[Jacobson 稠密性]\label{b2:thm:jacobson-density}
设 \(R\) 为含幺环、\(S\) 为简单左 \(R\) 模，\(D=\operatorname{End}_R(S)^{\mathrm{op}}\) 按 \(s\cdot f^{\mathrm{op}}=f(s)\) 作用于 \(S\) 的右侧。对任意右 \(D\) 线性无关的有限组 \(v_1,\ldots,v_n\)，以及任意 \(w_1,\ldots,w_n\in S\)，存在 \(r\in R\) 使
\[
rv_i=w_i\qquad(1\le i\le n).
\]
\end{theorem}

这称为\term[Jacobson-choumixing-dingli]{Jacobson 稠密性定理}[Jacobson density theorem]。存在的 \(r\) 通常不唯一；定理只在给定的有限组向量上指定其作用。

\begin{proof}
对 \(n\) 归纳。\(n=1\) 时，\(v_1\ne0\)，简单性给出 \(Rv_1=S\)，结论成立。设已对 \(n-1\) 成立。考虑左理想
\[
L=\{\,a\in R:av_1=\cdots=av_{n-1}=0\,\}.
\]
\(Lv_n\) 是 \(S\) 的子模，因此或为零，或为 \(S\)。先排除第一种情况。

若 \(Lv_n=0\)，归纳假设说明映射 \(a\mapsto(av_1,\ldots,av_{n-1})\) 从 \(R\) 满射到 \(S^{n-1}\)，因而可以定义
\[
F:S^{n-1}\longrightarrow S,
\qquad F(av_1,\ldots,av_{n-1})=av_n.
\]
若两个 \(a\) 给出相同输入，它们之差属于 \(L\)，所以输出之差也为零。这证明良定义。使用代表元的加法及左乘，得到 \(F\) 左 \(R\) 线性。令 \(f_i:S\to S\) 是 \(F\) 在第 \(i\) 个直和分量上的限制，则 \(f_i\in\Delta\)，且
\[
v_n=F(v_1,\ldots,v_{n-1})=\sum_{i<n}f_i(v_i)=\sum_{i<n}v_i d_i
\]
对某些 \(d_i\in D\) 成立，与给定向量组右线性无关矛盾。因此 \(Lv_n=S\)。

现在先用归纳假设找 \(r_0\) 使 \(r_0v_i=w_i\)、\(i<n\)。再取 \(\ell\in L\) 使 \(\ell v_n=w_n-r_0v_n\)。令 \(r=r_0+\ell\)，则前 \(n-1\) 个值因 \(\ell\in L\) 保持不变，最后一个值也达到要求，归纳完成。
\end{proof}

更一般地，给定任意右 \(D\) 线性算子 \(T\) 及有限组向量，先在它们的右线性生成空间中选基，再对这组基应用定理，即可找到 \(r\) 在原来的全部向量上与 \(T\) 一致。因此每个基本邻域 \(U(T;v_1,\ldots,v_m)\) 都与 \(\rho(R)\) 相交，即
\[
\overline{\rho(R)}^{\,\mathrm{fin}}=\operatorname{End}_D(S_D).
\]
有限拓扑中的稠密性只要求有限组上的同时一致；当 \(S_D\) 无限维时，它不要求某个 \(r\) 在整个空间上等于 \(T\)。

\subsection{有限维情形的矩阵代数}
\label{b2:subsec:1404:03}

\begin{corollary}\label{b2:cor:density-finite-dimensional}
设 \(R\) 为含幺环、\(S\) 为简单左 \(R\) 模，并令 \(D=\operatorname{End}_R(S)^{\mathrm{op}}\) 按 \(s\cdot f^{\mathrm{op}}=f(s)\) 作用于 \(S\) 的右侧。若 \(\dim_D S=n<\infty\)，则
\[
R/\operatorname{Ann}_R(S)\cong\operatorname{End}_D(S_D)\cong M_n(D).
\]
特别地，本原环若有一个对其自同态右除环有限维的忠实简单模，则它是单 Artin 环。
\end{corollary}
\begin{proof}
取右 \(D\) 基 \(v_1,\ldots,v_n\)。对任意 \(T\in\operatorname{End}_D(S)\)，在\cref{b2:thm:jacobson-density}中指定 \(w_i=T(v_i)\)，便找到 \(r\) 在全部基向量上与 \(T\) 相同。两者右线性，所以处处相同，作用同态满射。它的核为 \(\operatorname{Ann}_R(S)\)，由环同构定理得到第一个同构。

把向量写成右坐标列 \(\sum_i v_i d_i\)，若 \(T(v_j)=\sum_i v_i a_{ij}\)，则新坐标为 \((\sum_j a_{ij}d_j)_i\)。矩阵从左作用于右坐标列，复合对应通常的矩阵乘法，因此得到 \(M_n(D)\)，不再取反环。矩阵环的单性和 Artin 性由\cref{b2:prop:simple-artinian}及其矩阵块说明得到。
\end{proof}

这并没有证明每个本原环都为矩阵环；有限维条件是将有限组插值变成整个算子的等式所必需的。在前一节的无限维例子中，没有一个有限向量组能充当整个空间的基。

% !TeX root = ../../Book2.tex
% 纲要标题：### 14.5 Burnside 矩阵代数定理
% 纲要位置：计划纲要.md，第 1235 行。

\section{Burnside 矩阵代数定理}
\label{b2:sec:1405}

代数闭域上的有限维不可约作用具有更强的结论：自同态除环只剩底域标量，因而稠密性直接给出全部矩阵。这个结论通常称为 Burnside 矩阵代数定理；它对群作用也有用，但陈述的对象是线性算子组成的代数。

% 纲要要点（供目录核对）：
%
% * 代数闭域上有限维不可约矩阵代数
% * 与稠密性定理的联系
% * 区分轨道计数引理及有限群 p^a q^b 可解性定理
%
% ---
%

\subsection{代数闭域上的不可约矩阵代数}
\label{b2:subsec:1405:01}

\begin{theorem}[Burnside 矩阵代数]\label{b2:thm:burnside-matrix-algebra}
设 \(k\) 为代数闭域，\(V\ne0\) 有限维，\(A\subseteq\operatorname{End}_k(V)\) 为含恒等算子的 \(k\) 子代数。若 \(V\) 没有非零真 \(A\) 稳定子空间，则 \(A=\operatorname{End}_k(V)\)。
\end{theorem}
\begin{proof}
假设使 \(V\) 成为简单左 \(A\) 模。由\cref{b2:cor:schur-algebraically-closed}，\(\operatorname{End}_A(V)=k\operatorname{id}_V\)：有限维与代数闭性使每个自同态有一个特征值，Schur 引理再迫使该自同态等于相应标量。因此本节的右标量除环 \(D=\operatorname{End}_A(V)^{\mathrm{op}}\) 就是 \(k\)。

\cref{b2:cor:density-finite-dimensional}于是说明作用映射 \(A\to\operatorname{End}_k(V)\) 满射。而这个映射原本就是子代数的包含，故 \(A=\operatorname{End}_k(V)\)。
\end{proof}

这一\term[Burnside-juzhen-daishu-dingli]{Burnside 矩阵代数定理}[Burnside's matrix algebra theorem]没有假设 \(k\) 的特征为零。有限维保证特征多项式有定义且次数为正，代数闭性保证它有根，再由 Schur 引理迫使上述自同态为标量。

\ProseParagraphBreak
例如，对域 \(k\) 上的一族矩阵，若其生成的含幺代数作用不可约，且 \(k\) 代数闭，则这些矩阵的有限乘积与线性组合生成全部矩阵空间。结论是代数生成，不是说原先有限个矩阵本身已经线性生成了全部矩阵。

\subsection{与稠密性定理的联系}
\label{b2:subsec:1405:02}

一般稠密性使用右除环 \(D=\operatorname{End}_A(V)^{\mathrm{op}}\)。上一定理先把它识别为 \(k\)，然后在有限基上同时指定像。代数闭假设若删去，结果可以改变：复数左乘组成 \(M_2(\mathbb R)\) 中的子代数
\[
\left\{
\,\begin{pmatrix}a&-b\\b&a\end{pmatrix}
\ \middle|\ a,b\in\mathbb R\,
\right\}.
\]
它在实平面上不可约，因为任何非零真稳定子空间只能是一条实直线，却不可能对乘 \(i\) 的九十度旋转稳定。但这个代数只有二维，严格小于四维的 \(M_2(\mathbb R)\)。稠密性对此给出的目标是复线性自同态，恰好就是原代数。

\ProseParagraphBreak
有限维也不能直接删除。设 \(V\) 有可数基 \(e_1,e_2,\ldots\)，\(F\) 为全部有限秩算子组成的双边理想，取 \(A=k\operatorname{id}+F\)。它含有将任意非零向量送到任意目标向量的秩一算子，所以作用不可约。可是右移算子 \(T(e_i)=e_{i+1}\) 不属于 \(A\)：对每个 \(\lambda\in k\)，\(T-\lambda\operatorname{id}\) 都单射，因为有限支集向量的最后一个非零坐标在右移后产生一个无法被抵消的末项。因此其像无限维，不可能是有限秩算子。

\ProseParagraphBreak
这个例子仍能满足有限组插值：把有限组向量所生成的子空间扩充为直和分量，在该有限维部分按要求定义算子，在一个补空间上取零，就得到所需有限秩算子。因此“在任意指定的有限组上一致”与“作为整个空间上的算子相等”确实有不同的含义。

\ProseParagraphBreak

文献中 Burnside 的名字还用于有限群作用的轨道计数引理，以及“阶为 \(p^a q^b\) 的有限群可解”的定理。本节的矩阵代数定理与这两个陈述分别讨论不同对象，不能只凭同名互相引用。特别地，矩阵代数定理没有关于群阶只含两个素因子的条件，也不直接给出群的可解性。

\ProseParagraphBreak
本定理的有限维代数闭域版本，可对照 Etingof 等人的\cite[Section 3.2, Corollary 3.2.1, Theorem 3.2.2(i), p. 43]{EtingofEtAl2011RepresentationTheory}；该书将它列为稠密性定理的一部分。后续研究有限群表示时，群代数的像正是一类矩阵代数，因此这个结果会进入表示论；本章给出的证明只依赖 Schur 引理、自同态的特征值和\cref{b2:thm:jacobson-density}。


\begin{chaptersummary}
Jacobson 根既是全部极大左理想的交，也是全部极大右理想及简单左模零化子的交，所以是双边理想。\(x\in J(R)\) 的单位判据要求所有 \(1-rx\) 可逆；只检验 \(1-x\)，甚至只检验底域标量倍，都可能遗漏非交换环中的障碍。根在一般商下只保证像包含于商环的根，只有商掉的理想包含于原根时才得到等式。

\ProseParagraphBreak
Nakayama 引理把根的可逆性用于有限生成组的消减，继而给出商模生成元的提升。局部环上，最少生成元数由一个除环上的剩余维数计算；半单商有多个矩阵块时，则需分别比较简单因子重数与矩阵块大小。一般根理想下，有限生成性不能删除；若理想本身幂零，则可以反复代入直到理想的幂为零，不再需要有限生成。

\ProseParagraphBreak
半本原环可由一族本原商环的次直积描述，本原环则要求同一个简单模忠实。对简单左模 \(S\)，置 \(D=\operatorname{End}_R(S)^{\mathrm{op}}\)。稠密性定理及其关系判据说明，有限组指定像能够由同一个环元素实现，当且仅当它们保留原向量间的全部右 \(D\) 线性关系。因此作用像在有限拓扑中的闭包为 \(\operatorname{End}_D(S_D)\)，这个闭包未必等于作用像本身。只有当该右除环上的维数有限时，有限组条件才覆盖整个算子；Burnside 定理再利用代数闭性把右除环缩为底域，从而得到全部矩阵。下一章将把 Jacobson 根与有限维性结合，研究根滤过和投射覆盖。
\end{chaptersummary}

\BookExercises

\begin{exercise}\label{b2:ex:14-scalar-unit-test}
设 \(k\) 为任意域。给出 \(x\in M_2(k)\)，使 \(1-\lambda x\) 对每个 \(\lambda\in k\) 都可逆，但 \(x\notin J\bigl(M_2(k)\bigr)\)。明确指出全环单位判据怎样排除这个 \(x\)。
\end{exercise}
\begin{solution}
取 \(x=E_{12}\)，则 \(x^2=0\)，所以 \((1-\lambda x)^{-1}=1+\lambda x\) 对任意标量成立。但取环元素 \(r=E_{21}\)，有 \(rx=E_{22}\)，故 \(1-rx=\operatorname{diag}(1,0)\) 不可逆。由\cref{b2:thm:jacobson-unit-criterion}，\(x\notin J\bigl(M_2(k)\bigr)\)。标量倍只检验了全环中很小的一部分元素，不能替代对全部 \(r\) 的条件。
\end{solution}

\begin{exercise}\label{b2:ex:14-integer-quotient-radical}
设 \(n=\prod_{i=1}^s p_i^{a_i}\ge2\)，其中 \(p_i\) 是两两不同的素数，\(a_i\) 为正整数。求 \(J(\mathbb Z/n\mathbb Z)\)，并求使其幂为零的最小正整数。
\end{exercise}
\begin{solution}
整数商环的极大理想是 \((\bar p_i)\)：其商为 \(\mathbb F_{p_i}\)，反向每个极大理想的原像是包含 \(n\mathbb Z\) 的整数极大理想，故只能是这些素数理想。因此根由
\[
\overline{p_1\cdots p_s}
\]
生成。其 \(r\) 次幂为零当且仅当 \(n\) 整除 \((p_1\cdots p_s)^r\)，即每个 \(r\ge a_i\)。最小正整数为 \(\max_i a_i\)。尤其 \(n\) 无平方因子时，根为零，此时最小正整数为一。
\end{solution}

\begin{exercise}\label{b2:ex:14-product-radical}
设 \((R_i)_{i\in I}\) 为由集合 \(I\) 编号的含幺环族。证明其直积满足
\[
J\left(\prod_{i\in I}R_i\right)=\prod_{i\in I}J(R_i).
\]
论证应同时适用于无限指标集。
\end{exercise}
\begin{solution}
直积中的元素可逆，当且仅当每个坐标可逆：逆元只能逐坐标取逆，而这些逆确实组成直积元素。若 \(x_i\in J(R_i)\)，则对任意 \(r=(r_i)\)，各 \(1-r_ix_i\) 都可逆，故 \(1-rx\) 可逆，得到正向包含。反过来，若 \(x\) 属于直积环的根，固定 \(i\) 及 \(a\in R_i\)，令 \(r\) 仅第 \(i\) 坐标为 \(a\)，其余为零。\(1-rx\) 可逆迫使 \(1-ax_i\) 可逆。对任意 \(a\) 使用单位判据，得 \(x_i\in J(R_i)\)，证明反向包含。没有一步要求元素只有有限支集。
\end{solution}

\begin{exercise}\label{b2:ex:14-matrix-root-powers}
设 \(R\) 为含幺环，\(I\) 为其双边理想，\(n,r\) 为正整数。证明 \(M_n(I)^r=M_n(I^r)\)。据此计算 \(J\bigl(M_3(\mathbb Z/12\mathbb Z)\bigr)\) 及其幂零指数。
\end{exercise}
\begin{solution}
矩阵乘积的条目是 \(r\) 个 \(I\) 元素乘积的有限和，故有包含 \(M_n(I)^r\subseteq M_n(I^r)\)。反向，\(I^r\) 中的每个元素为 \(a_1\cdots a_r\) 的有限和，且
\[
(a_1E_{i1})\left(\prod_{\nu=2}^{r-1}(a_\nu E_{11})\right)(a_rE_{1j})
=a_1\cdots a_r E_{ij}
\]
在 \(r\ge2\) 时成立，其中 \(r=2\) 时中间的空积取单位矩阵；\(r=1\) 直接成立。把任意矩阵拆成这些矩阵单位位置的和，得到反向包含。

\(J(\mathbb Z/12\mathbb Z)=(\bar6)\)，非零且平方为零。由\cref{b2:prop:matrix-radical}，所求根为全部条目在 \((\bar6)\) 中的三阶矩阵。上式说明其平方为零，而 \(\bar6E_{11}\ne0\)，所以幂零指数为二。
\end{solution}

\begin{exercise}\label{b2:ex:14-radical-general-map}
根在满同态下的像包含于目标的根。给出一个非满的保单位元环同态，使这个包含失败。可以使用局部环 \(k[t]_{(t)}\) 到其分式域的包含。
\end{exercise}
\begin{solution}
令 \(R=k[t]_{(t)}\subseteq k(t)\)。在 \(R\) 中，分母在零点非零；一个分式可逆当且仅当分子也在零点非零。因此非单位为理想 \((t)\)，它是唯一极大理想，\(J(R)=(t)\)。目标 \(k(t)\) 为域，根为零；包含映射却把 \(t\in J(R)\) 送到非零元素 \(t\)。所以包含失败。满同态的证明依赖目标简单模限制后仍简单；一般包含映射没有这一保证。
\end{solution}

\begin{exercise}\label{b2:ex:14-units-lift}
设 \(k=\mathbb F_q\)，\(R=k[\varepsilon]/(\varepsilon^3)\)，以下仍用 \(\varepsilon\) 表示其商类。求 \(|R^\times|\) 与 \(|\operatorname{GL}_2(R)|\)，并写出矩阵
\[
A=\begin{pmatrix}1&\varepsilon\\\varepsilon&1+\varepsilon\end{pmatrix}
\]
的逆。说明从 \(\operatorname{GL}_2(R)\) 到 \(\operatorname{GL}_2(k)\) 的约化映射的核。
\end{exercise}
\begin{solution}
理想 \(I=(\varepsilon)\) 满足 \(I^3=0\)，故包含于 \(J(R)\)。由\cref{b2:cor:units-lift-radical}，元素为单位当且仅当其常数项非零，所以 \(|R^\times|=(q-1)q^2\)。同一推论使矩阵约化映射满射：任意可逆的 \(k\) 矩阵任选逐项提升都可逆。其核恰为
\[
1+M_2(I)=\{\,I_2+B\mid B\in M_2(I)\,\}.
\]
每个这样的矩阵确实可逆，因为 \(B^3=0\)。\(I\) 有 \(q^2\) 个元素，四个条目可独立选取，故核有 \(q^8\) 个元素。因此
\[
|\operatorname{GL}_2(R)|=q^8(q^2-1)(q^2-q),
\]
其中两因子分别数出非零第一列及不在其线性生成空间中的第二列。

写 \(A=I_2+\varepsilon C\)，其中 \(C=\begin{pmatrix}0&1\\1&1\end{pmatrix}\)。由 \(\varepsilon^3=0\)，
\[
A^{-1}=I_2-\varepsilon C+\varepsilon^2C^2
=\begin{pmatrix}
1+\varepsilon^2&-\varepsilon+\varepsilon^2\\
-\varepsilon+\varepsilon^2&1-\varepsilon+2\varepsilon^2
\end{pmatrix}.
\]
这个有限几何级数公式在特征二时也成立，整数系数按 \(k\) 的特征解释。
\end{solution}

\begin{exercise}\label{b2:ex:14-triangular-radical}
设 \(k\) 为域，考虑分块上三角矩阵环
\[
R=\left\{\,
\begin{pmatrix}A&B\\0&C\end{pmatrix}
\ \middle|\ A\in M_2(k),\ B\in M_{2\times3}(k),\ C\in M_3(k)
\,\right\}.
\]
求 \(J(R)\)、其幂零指数及 \(R/J(R)\)。比较它与通常的 \(T_5(k)\)：为什么矩阵总大小都是五，根的幂零指数却不同？
\end{exercise}
\begin{solution}
取 \(N\) 为两个对角块都为零的矩阵集合。它是双边理想，且 \(N^2=0\)。由\cref{b2:cor:nil-ideal-radical}，\(N\subseteq J(R)\)。对角块映射给出
\[
R/N\cong M_2(k)\times M_3(k).
\]
两块的自然列模分别简单且忠实。把它们看成直积环的两个简单模，零化子之交为零，故 \(J(R/N)=0\)。由\cref{b2:prop:radical-quotient}，\(J(R)/N=0\)，所以 \(J(R)=N\)。\(N\ne0\) 且 \(N^2=0\)，其幂零指数为二。

\cref{b2:prop:triangular-radical}给出 \(T_5(k)\) 的根幂零指数为五。它的严格上三角根中有从第一位置经第二、三、四到第五位置的四次非零乘积；在这里的两块环中，根只连接第二块到第一块，两个根元素的乘积已经为零。指数取决于根中允许的乘积，不由矩阵的总大小单独决定。
\end{solution}

\begin{exercise}\label{b2:ex:14-nakayama-radical-needed}
取 \(R=k\times k\)、\(I=k\times0\)、\(M=I\)，均采用坐标乘法。证明 \(M\) 为非零有限生成模且 \(IM=M\)。说明这为何不与 Nakayama 引理矛盾。
\end{exercise}
\begin{solution}
\(M\) 由 \((1,0)\) 生成，故非零且循环；\(I^2=I\)，所以 \(IM=M\)。但 \(R\) 的两个简单坐标模的零化子交为零，故 \(J(R)=0\)，非零理想 \(I\) 不包含于根。也可以直接取 \(x=(1,0)\in I\)，此时 \(1-x=(0,1)\) 不可逆，单位判据排除 \(x\in J(R)\)。Nakayama 引理的根理想假设未满足。
\end{solution}

\begin{exercise}\label{b2:ex:14-nakayama-finite-needed}
设 \(R=k[t]_{(t)}\)，\(M=k(t)/R\)。证明 \(M\ne0\)、\(tM=M\)，且每个元素都被某个 \(t\) 的幂零化。证明 \(M\) 不有限生成，并说明 \(M/J(R)M\) 的值。
\end{exercise}
\begin{solution}
\(t^{-1}\notin R\)，所以 \(M\ne0\)。任意 \(f+R\) 都是 \(t(f/t+R)\)，故 \(tM=M\)。把有理函数的分母写成 \(t^n\cdot q(t)\)、\(q(0)\ne0\)，便有 \(t^nf\in R\)，说明每个陪集被某个 \(t^n\) 零化。

若有限个陪集生成 \(M\)，取共同指数 \(N\) 零化这些生成元，交换性使 \(t^N\) 零化它们的全部 \(R\) 线性组合。但 \(t^{-N-1}+R\) 乘 \(t^N\) 后为非零的 \(t^{-1}+R\)，矛盾。所以不有限生成。\(J(R)=(t)\)，因此 \(J(R)M=M\)，剩余模为零；在缺少有限生成时，零剩余模并不能检测原模为零。
\end{solution}

\begin{exercise}\label{b2:ex:14-nilpotent-ideal-generators}
设 \(R\) 为含幺环、\(I\) 为满足 \(I^r=0\) 的双边理想，\(r\ge1\)。对任意幺性左 \(R\) 模 \(M\)，证明 \(M\) 有限生成当且仅当 \(M/IM\) 在 \(R/I\) 上有限生成；在此情形证明
\[
\mu_R(M)=\mu_{R/I}(M/IM).
\]
若只假设 \(I\subseteq J(R)\)，上述有限生成性反向蕴含未必成立。用 \(R=k[t]_{(t)}\)、\(I=(t)\)、\(M=R\oplus k(t)\) 验证这一点。
\end{exercise}
\begin{solution}
有限生成模的商仍有限生成。反过来，提升 \(M/IM\) 的有限生成组为 \(x_1,\ldots,x_n\)，记 \(N=\sum_iRx_i\)，则 \(M=N+IM\)。由\cref{b2:cor:nakayama-generators}的幂零理想版本，不需预设 \(M\) 有限生成便有 \(M=N\)。任意生成组取像给出商模的生成组，任意商模生成组提升又给出原模生成组，两方向比较元素数便得最少数相等；零模也包括在内。

在题给局部环中，\(J(R)=(t)\)，且 \(tk(t)=k(t)\)。所以 \(M/IM\cong R/(t)\cong k\)，只需一个生成元；但 \(M\) 不有限生成，因为它的商模 \(k(t)\) 不有限生成。后者在零点的极点阶无统一上界，而有限个有理函数的 \(R\) 线性组合的极点阶有共同上界。这里 \(I\) 的任何幂均不为零，幂零理想下的反向提升论证因而不能使用。
\end{solution}

\begin{exercise}\label{b2:ex:14-local-generator-computation}
设 \(k\) 为域，令 \(R=k[x,y]/(x,y)^3\)，仍用 \(x,y\) 表示其商类，记 \(\mathfrak m=(x,y)\)。确定 \(J(R)\)，计算模 \(\mathfrak m\)、\(\mathfrak m^2\)、\(\mathfrak m^2\oplus R/(x)\) 的最少生成元个数。说明 \(x+y^2,y+x^2\) 是否生成 \(\mathfrak m\)。
\end{exercise}
\begin{solution}
\(\mathfrak m^3=0\)，且 \(R/\mathfrak m=k\)，故由幂零理想包含于根及商环根公式，\(J(R)=\mathfrak m\)。\(\mathfrak m/\mathfrak m^2\) 有基 \(x,y\)，所以 \(\mu_R(\mathfrak m)=2\)。\(\mathfrak m\mathfrak m^2=0\)，而 \(\mathfrak m^2\) 有基 \(x^2,xy,y^2\)，所以 \(\mu_R(\mathfrak m^2)=3\)。最后一个模的剩余空间由前三个基元素再加 \(R/(x)\) 的常数类组成，维数为四，故最少需要四个生成元。

\(x+y^2,y+x^2\) 在 \(\mathfrak m/\mathfrak m^2\) 中的像仍为 \(x,y\)，是一组基。目标模 \(\mathfrak m\) 有限生成，由\cref{b2:cor:nakayama-generators}，这两个提升生成 \(\mathfrak m\)，并且是不可删减生成组。
\end{solution}

\begin{exercise}\label{b2:ex:14-semisimple-generator-packing}
设 \(k\) 为域，令 \(R=M_2(k)\times M_3(k)\)，以 \(S_1=k^2,S_2=k^3\) 表示相应的简单列模。求 \(M=S_1^5\oplus S_2^7\) 的最少生成元个数，并构造达到这个数的生成组。
\end{exercise}
\begin{solution}
由\cref{b2:prop:semisimple-top-generators}，
\[
\mu_R(M)=\max\{\lceil5/2\rceil,\lceil7/3\rceil\}=3.
\]
构造三个向量 \(g_1,g_2,g_3\)。在五个 \(S_1\) 坐标中，\(g_j\) 的第 \(2j-1,2j\) 个位置分别放标准向量 \(e_1,e_2\)，超过五的位置删去，其余位置为零。在七个 \(S_2\) 坐标中，\(g_j\) 的第 \(3j-2,3j-1,3j\) 个位置分别放 \(e_1,e_2,e_3\)，超过七的位置删去，其余为零。

对每个 \(g_j\)，二阶矩阵的两列可以独立指定其负责的两个 \(S_1\) 坐标，三阶矩阵的三列可以独立指定其负责的三个 \(S_2\) 坐标；直积环中这两个矩阵又可独立选取。各 \(j\) 负责的坐标不交，合起来覆盖全模，所以三者生成 \(M\)。
\end{solution}

\begin{exercise}\label{b2:ex:14-irredundant-not-minimum}
设 \(R=k^4\)，采用坐标乘法，\(g_1,\ldots,g_m\in R\)。用各 \(g_j\) 的非零坐标位置给出它们生成左正则模 \({}_RR\) 的判据，以及这组生成元不可删减的判据。据此构造不可删减的三元生成组，证明任一不可删减有限生成组至多含四个元素，并求 \(\mu_R(R)\)。
\end{exercise}
\begin{solution}
记 \(A_j=\{\,i\in\{1,2,3,4\}\mid(g_j)_i\ne0\,\}\)。每个系数的四个坐标可独立选取，因此这组元素生成 \(R\) 当且仅当 \(\bigcup_jA_j=\{1,2,3,4\}\)：若漏掉某个位置，所有线性组合在那里为零；若每个位置均覆盖，利用坐标幂等元及域中逆元便可得到各个标准坐标向量。

在生成条件成立时，删去 \(g_j\) 后不再生成，当且仅当 \(A_j\) 含有一个不属于任何其他 \(A_\ell\) 的位置。因此不可删减等价于每个生成元都有这样的独占坐标。可取
\[
g_1=(1,1,0,0),\qquad g_2=(0,1,1,0),\qquad g_3=(0,0,0,1).
\]
它们分别有独占坐标一、三、四，故不可删减。给每个不可删减生成元选择一个独占坐标，这些坐标两两不同，故 \(m\le4\)。但单位元 \((1,1,1,1)\) 单独生成 \(R\)，而 \(R\ne0\)，所以 \(\mu_R(R)=1\)。
\end{solution}

\begin{exercise}\label{b2:ex:14-polynomial-semiprimitive}
对任意域 \(k\)，证明 \(J\bigl(k[t]\bigr)=0\)，但 \(k[t]\) 既不本原，也不半单。不使用“极大理想的交显然为零”作为证明。
\end{exercise}
\begin{solution}
对任意非零多项式 \(f\)，取 \(r=t\)，则 \(1-tf\) 具有正次数，不是 \(k[t]\) 的单位。由\cref{b2:thm:jacobson-unit-criterion}，\(f\notin J\bigl(k[t]\bigr)\)，所以根为零。交换本原环必须为域，而 \(t\) 不可逆，所以不本原。降链 \((t)\supsetneq(t^2)\supsetneq\cdots\) 不稳定，故不是 Artin 环；半单环由\cref{b2:thm:artin-wedderburn}具有 Artin 性，所以它不半单。
\end{solution}

\begin{exercise}\label{b2:ex:14-infinite-product}
设 \(k\) 为域，\(R=\prod_{n\ge1}k\)。证明每个非零元素都能被某个简单左模检测到，因而 \(R\) 半本原；再证明它没有忠实简单模，也不是半单环。
\end{exercise}
\begin{solution}
第 \(n\) 个坐标投影使 \(k\) 成为简单 \(R\) 模。如果 \(r\ne0\)，某个坐标 \(r_n\ne0\)，它在该简单模上作用非零。于是全部简单模零化子的交为零，由\cref{b2:prop:radical-simple-annihilators}，\(J(R)=0\)。\(R\) 交换且有非平凡幂等元，不是域，故由\cref{b2:prop:commutative-primitive-field}不本原，没有忠实简单模。令 \(I_n\) 为前 \(n\) 个坐标都为零的理想，则 \(R\supsetneq I_1\supsetneq I_2\supsetneq\cdots\)，故不 Artin，也不半单。
\end{solution}

\begin{exercise}\label{b2:ex:14-evaluation-subdirect}
设 \(k\) 为无限域，考虑取值映射 \(k[t]\rightarrow\prod_{a\in k}k\)、\(f\mapsto\bigl(f(a)\bigr)_a\)。证明它给出一个严格小于全直积的次直积。若 \(k\) 为有限域，这个映射的核与像又是什么？
\end{exercise}
\begin{solution}
无限域上，一个非零多项式只有有限多个根，所以取值映射单射。任一坐标映射满射，因为常数多项式可以取任意值。它不到全直积满射：只在 \(a=0\) 处取一、在其他点取零的坐标族不可能来自多项式，否则这个多项式有无限多个根，须为零，又与在零点的值矛盾。

有限域上，令 \(P(t)=\prod_{a\in k}(t-a)\)。一个多项式在每点取零，当且仅当逐个一次因子都整除它；这些因子两两互素，故等价于 \(P\mid f\)，核为 \((P)\)。对每个 \(a\)，Lagrange 多项式
\[
\ell_a(t)=\prod_{\substack{b\in k\\b\ne a}}\frac{t-b}{a-b}
\]
在 \(a\) 处取一，其余点取零，所以任意坐标族由有限和 \(\sum_a c_a\ell_a\) 给出，映射满射。于是 \(k[t]/(P)\cong\prod_{a\in k}k\)，有限与无限乘积的结论不同。
\end{solution}

\begin{exercise}\label{b2:ex:14-right-division-scalars}
设 \(D_0\) 为任意除环，\(R=D_0\)、\(S={}_{D_0}D_0\)。计算 \(\Delta=\operatorname{End}_R(S)\) 及本章采用的右标量除环 \(D=\Delta^{\mathrm{op}}\)。对 \(0\ne v\in S,w\in S\)，写出实现 \(rv=w\) 的 \(r\)，并说明为何 \(S_D\) 的无关组最多只有一个向量。
\end{exercise}
\begin{solution}
每个左线性自同态由 \(f(1)=a\) 决定，并为右乘 \(f(x)=xa\)。两个右乘的复合反转参数乘法，所以 \(\Delta\cong D_0^{\mathrm{op}}\)，再取反环得到 \(D\cong D_0\)。按此识别，右标量作用就是通常的右乘，\(S_D\) 以一为基，右维数为一。

取 \(r=wv^{-1}\)，有 \(rv=w\)。任意两个非零向量满足 \(v_2=v_1(v_1^{-1}v_2)\)，所以右线性相关。即使 \(D_0\) 在某个中心子域上有更高维数，也不能对这种右相关组任意指定像；正确的插值条件使用 \(D\)。
\end{solution}

\begin{exercise}\label{b2:ex:14-density-relations}
设 \(D_0\) 为除环、\(R=M_2(D_0)\)、\(S=D_0^2\) 为简单列模，右标量作用采用通常右乘。固定 \(a,b\in D_0\)，令
\[
v_1=e_1,\qquad v_2=e_2,\qquad v_3=e_1a+e_2b.
\]
给定 \(w_1,w_2,w_3\in S\)，求同时满足 \(rv_i=w_i\) 的全部 \(r\in R\)。若只给定 \(rv_1=w_1\) 与 \(rv_3=w_3\)，分别讨论 \(b=0\) 和 \(b\ne0\)，并注意非交换系数的位置。
\end{exercise}
\begin{solution}
与全部矩阵单位及 \(D_0\) 条目作用交换的自同态，恰好在两个坐标上分别右乘同一个 \(D_0\) 元素，所以 \(\operatorname{End}_R(S)\cong D_0^{\mathrm{op}}\)，相应右标量除环确为 \(D_0\)。同时指定前两列已唯一确定矩阵 \(r=(w_1\ w_2)\)。右线性要求
\[
w_3=w_1a+w_2b,
\]
这也充分：满足时上述矩阵实现全部指定像。原向量的任意关系均为 \((d_1,d_2,d_3)=(-ad,-bd,d)\)，\(d\in D_0\)，所以这一条件也正是\cref{b2:cor:density-compatible-relations}的全部关系保持条件。

若只指定第一、第三个像，当 \(b=0\) 时必须且只需 \(w_3=w_1a\)，第二列可任取。当 \(b\ne0\) 时，第一、第三个向量右线性无关；第二列必须为
\[
(w_3-w_1a)b^{-1},
\]
故任意 \(w_1,w_3\) 都对应唯一矩阵。因为这里是右标量关系，\(b^{-1}\) 乘在向量的右侧；把它移到左侧在一般除环上既没有相同含义，也没有相同结果。
\end{solution}

\begin{exercise}\label{b2:ex:14-finite-rank-commutant}
设 \(V=\bigoplus_{i\ge1}ke_i\)，\(F\) 为全部有限秩算子，\(F_0\) 为相对于此基只有有限多个非零矩阵条目的算子。证明 \(F_0\) 是不含恒等算子的 \(k\) 子代数，且 \(F_0\subsetneq F\)。证明 \(F_0\) 在 \(\operatorname{End}_k(V)\) 的有限拓扑中已经稠密，并比较 \(k\operatorname{id}+F_0\) 与 \(k\operatorname{id}+F\)。
\end{exercise}
\begin{solution}
有限多个非零条目的矩阵相加、数乘和相乘仍只有有限多个非零条目，所以 \(F_0\) 为子代数；每个这样的算子像均有限维，故属于 \(F\)。恒等矩阵有无限多个非零对角项，不在 \(F_0\)。算子 \(U(e_i)=e_1\) 的像为 \(ke_1\)，但有无限多个非零条目，说明 \(F_0\subsetneq F\)。

给定目标算子 \(T\) 及有限个测试向量 \(v_1,\ldots,v_n\)，它们总共只使用有限个基位置，记其集合为 \(J\)。令 \(B(e_j)=T(e_j)\) 对 \(j\in J\) 成立，在其余基向量上取零。有限多个 \(T(e_j)\) 又各有有限支集，所以 \(B\) 的非零矩阵条目只有有限多个，\(B\in F_0\)；并且 \(B(v_i)=T(v_i)\)。于是每个基本邻域都与 \(F_0\) 相交，其闭包为全部自同态环。

有 \(k\operatorname{id}+F_0\subsetneq k\operatorname{id}+F\)：上面的 \(U\) 不能写成 \(c\operatorname{id}+B\)、\(B\in F_0\)。若 \(c=0\)，\(U\) 有无限多个非零条目；若 \(c\ne0\)，\(U-c\operatorname{id}\) 在所有 \(i\ge2\) 的对角位置都非零，同样不在 \(F_0\)。由\cref{b2:prop:finite-rank-density-counterexample}，后者还严格小于 \(\operatorname{End}_k(V)\)。这两个含幺真子代数却都有同样的有限拓扑闭包，说明稠密性不能区别它们在无限多基位置上的限制。
\end{solution}

\begin{exercise}\label{b2:ex:14-burnside-rational-counterexample}
令 \(J=\begin{pmatrix}0&2\\1&0\end{pmatrix}\)，\(A=\mathbb Q[J]\subseteq M_2(\mathbb Q)\)。证明 \(A\) 在 \(\mathbb Q^2\) 上作用不可约，但不等于全部矩阵环；求其中心化子 \(\operatorname{End}_A(\mathbb Q^2)\)，并说明稠密性给出的正确目标。
\end{exercise}
\begin{solution}
\(J^2=2I\)，而 \(t^2-2\) 在 \(\mathbb Q\) 上不可约，所以 \(A\cong\mathbb Q(\sqrt2)\)。把向量 \((a,b)\) 识别为 \(a+b\sqrt2\)，\(J\) 对应乘 \(\sqrt2\)。任意非零 \(A\) 稳定子空间都是这个域的非零左理想，故为全体，作用不可约。\(A\) 的有理维数为二，严格小于矩阵环的四。

令 \(C=\begin{pmatrix}u&v\\w&z\end{pmatrix}\)，等式 \(CJ=JC\) 给出 \(u=z,v=2w\)，所以 \(C=uI+wJ\in A\)。因此中心化子恰为 \(A\)，它是交换域；相应右标量除环也是 \(A\)，\(\mathbb Q^2\) 在它上面一维，稠密性给出的全部右线性自同态正是 \(A\)。缺少代数闭性时，不能把这个目标改成全部有理线性自同态。
\end{solution}

\begin{exercise}\label{b2:ex:14-algebra-generation}
设 \(k\) 为任意域，令 \(U=I+E_{12}\)、\(V=I+E_{21}\)。证明 \(U,V\) 生成的含幺 \(k\) 子代数是 \(M_2(k)\)，却不能仅用 \(I,U,V\) 的线性组合得到所有矩阵。说明这里的结论为什么不需要 \(k\) 代数闭。
\end{exercise}
\begin{solution}
由 \(U-I=E_{12}\)、\(V-I=E_{21}\)，生成代数包含这两个矩阵单位，继而包含乘积 \(E_{12}E_{21}=E_{11}\)、\(E_{21}E_{12}=E_{22}\)。四个矩阵单位是 \(M_2(k)\) 的基，所以生成代数为全体。\(I,U,V\) 的线性生成空间只有三维：它们与 \(I,E_{12},E_{21}\) 具有相同线性生成空间，后者显然无关。乘积提供了原线性组合中缺少的方向。这里是显式矩阵运算，对任意域成立；Burnside 定理的代数闭条件是其一般不可约性判据的条件，并非所有具体生成等式成立的必要条件。
\end{solution}

\begin{exercise}\label{b2:ex:14-group-degree-bound}
设有限群 \(G\) 在代数闭域 \(k\) 上有非零有限维不可约表示 \(\rho:G\rightarrow\operatorname{GL}(V)\)，记 \(d=\dim_kV\)。证明 \(\rho(G)\) 的线性生成空间为 \(\operatorname{End}_k(V)\)，并推出 \(d^2\le |G|\)。不要加入关于特征不整除群阶的假设。
\end{exercise}
\begin{solution}
记 \(A=\operatorname{span}_k\rho(G)\)。群乘法使 \(\rho(g)\rho(h)=\rho(gh)\)，所以 \(A\) 对乘法封闭，且包含恒等算子。子空间对全部群元素稳定，当且仅当对其线性生成空间稳定，故 \(A\) 的作用不可约。\cref{b2:thm:burnside-matrix-algebra}于是给出 \(A=\operatorname{End}_k(V)\)。左侧由至多 \(|G|\) 个矩阵生成，右侧维数为 \(d^2\)，所以 \(d^2\le|G|\)。这个不等式对任意特征成立。
\end{solution}
